% Part: proof-theory
% Chapter: cut-elimination
% Section: introduction

\documentclass[../../../include/open-logic-section]{subfiles}

\begin{document}

\olfileid{pt}{cut}{int}

\olsection{مقدمة}

قاعدة~\Cut{} هي
\[
\Axiom$\Gamma \fCenter \Delta, !A$
\Axiom$!A, \Pi \fCenter \Lambda$
\RightLabel{\Cut}
\BinaryInf$\Gamma, \Pi \fCenter \Delta, \Lambda$
\DisplayProof
\]
(وتسمى~!!{formula}~$!A$ \emph{صيغة القطع}.)

أدرج غنتسن قاعدة~\Cut{} في حسابه الأصلي للمتتاليات~\Log{LK}. وتنص
\emph{مبرهنة غنتسن الرئيسة} (\emph{Hauptsatz}) على أن كل متتالية
!!{provable} في~\Log{LK} هي أيضًا~!!{provable} من دون استخدام قاعدة
\Cut{}؛ أي يمكن «حذف»~\Cut{} من براهين~\Log{LK}. ولا يحتوي نظامانا
\Log{G1c} و\Log{G3c} قاعدة~\Cut، لكن يمكن طرح السؤال نفسه عنهما: هل
يمكن حذف~\Cut{} من~!!{proof}s التي تستخدم~\Cut{} إلى جانب سائر قواعد
\Log{G1c} (أو قواعد~\Log{G3c})؟ وبمصطلحاتنا المعتمدة، يكافئ ذلك السؤال:
هل~\Cut{} قاعدة مقبولة في~\Log{G1c} (أو~\Log{G3c})؟

لعل مبرهنة حذف القطع أهم نتائج نظرية البرهان وأشهرها. فهي تثبت أولًا أن
\Log{G1c} و\Log{G3c} لا يحتاجان إلى قاعدة قد تبدو، للوهلة الأولى،
ضرورية بوضوح. فعندما تحاول صورنة براهين فعلية، تلاحظ أنك تحتاج إلى
طريقة لمعالجة استخدام اللمم. يحب الرياضيون أن يبرهنوا نتائج ثم يستندوا
إليها في براهين نتائج أخرى. فقد تصوغ أولًا برهان النتيجة~$L$ (اللمّة)
بوصفه~!!{proof} للمتتالية~$\Sequent L$. ثم تصوغ برهان النتيجة الثانية
$R$، الذي يستخدم~$L$، بوصفه~!!{proof} للمتتالية~$L \Sequent R$. فكيف
تعرف أن لـ$R$ برهانًا، أي أن للمتتالية~$\Sequent R$~!!a{proof}؟ تحتاج
إلى طريقة لضم~!!{proof}s الاثنين، وتتيح لك قاعدة~\Cut{} ذلك.

ذكرنا أن~\Log{G1c} و\Log{G3c} تامان، ومن ثم يبرهنان كل متتالية تتوقع
إمكان برهنتها، أي كل المتتاليات الصادقة. ويمكن إثبات ذلك مباشرة. لكن
حين كان غنتسن يكتب، لم يكن التمام معروفًا إلا لأنظمة الاشتقاق
المسلّمية. ولإثبات تمام~\Log{LK}، قدم غنتسن ترجمة لـ!!{derivation}s في
النظام المسلّمي إلى~!!{proof}s في~\Log{LK}. وكانت هذه الترجمة تستخدم
\Cut{} استخدامًا جوهريًا. وليس حذف القطع نتيجة مهمة للمنطق الكلاسيكي
وحده؛ فلمنطقات كثيرة حسابات متتاليات. وفي بعض المنطقات يكون البرهان
المباشر على تمام حساب متتالياتها صعبًا أو مستحيلًا، فتكون الترجمات من
أنظمة أخرى، باستخدام~\Cut، السبيل الوحيد لإثبات التمام. وعندئذ يلزم
إثبات حذف القطع برهانيًا لبيان أن~\Cut{} زائدة وأن النظام الخالي منها
تام.

ويهم حذف القطع أيضًا لأنه يفضي إلى نتائج أخرى ذات أهمية مستقلة.
وسنتناول مثالين هما لمّة كريغ للاستيفاء ومبرهنة هيربراند.

تقوم الأفكار المستخدمة عادة في برهان حذف القطع على بيان كيفية تحويل
!!a{proof} يستخدم~\Cut{} إلى برهان لا يستخدمها. وغالبًا ما تكون هذه
التحويلات «موضعية»: نأخذ جزءًا من~!!a{proof}، يكون عادة~!!{proof}
جزئيًا ينتهي باستدلال~\Cut{}، ونستبدل به~!!{proof} جزئيًا آخر للمتتالية
نفسها حُذف منه استدلال~\Cut{} أو استُبدل باستدلالات أخرى. وتوفر هذه
الحجج خوارزميات خطوة بخطوة لتحويل~!!{proof}s التي تستخدم~\Cut{} إلى
براهين لا تستخدمها، فتقدم معلومات أكثر حين تُستخدم إجراءات حذف القطع
في التطبيقات. فمثلًا، ثمة براهين نموذجية نظرية للمّة الاستيفاء تأخذ
الصورة: «إذا كانت~$!A \lif !B$ صادقة، وُجدت صيغة استيفاء» (ولا يهم
الآن معنى صيغة الاستيفاء). أما الحجة البرهانية التي تستخدم حذف القطع
فتقدم خوارزمية تأخذ~!!a{proof} لـ$!A \lif !B$ وتعيد صيغة استيفاء.
وخلافًا للحجة النموذجية النظرية، تكون الحجة البرهانية
\emph{بنائية}.

ثمة طريقتان شائعتان لبرهنة مبرهنة حذف القطع. تبين إحداهما، التي ترجع
إلى غنتسن، إمكان إزالة القطوع العليا في البرهان، ثم تزيل القطوع العليا
تباعًا حتى لا يبقى شيء منها. أما الأخرى، التي ترجع أصلًا إلى شوتّه
وتايت، فتركز على أعقد القطوع في~!!a{proof}، وتزيل تباعًا القطوع القصوى
بمعنى أنه لا يوجد قطع آخر له~!!{formula} قطع أعمق. ولكل منهما مزايا
وعيوب. وقد ينجح أحد النهجين في نظام ما بينما يصعب تنفيذ الآخر أو
يستحيل. لذلك نصف النهجين كليهما. وسنبرهن حذف القطع لـ\Log{G3c}.
والواقع أننا لن نبين إمكان حذف~\Cut{}، بل إمكان حذف
\emph{القطع المشترك السياق}~\CutCS:
\[
\Axiom$\Gamma \fCenter \Delta, !A$
\Axiom$!A, \Gamma \fCenter \Delta$
\RightLabel{\CutCS}
\BinaryInf$\Gamma \fCenter \Delta$
\DisplayProof
\]
إذا كان حساب المتتاليات يجيز الإضعاف والتقليص، سواء بوصفهما قاعدتين
فعليتين كما في~\Log{G1c} أو قاعدتين مقبولتين كما في~\Log{G3c}، أمكن
لـ\Cut{} محاكاة~\CutCS والعكس. ونختار~\CutCS{} بدلًا من~\Cut لأن
استراتيجية حذف القطع الأولى أسهل وصفًا لـ\CutCS، ولأن الاستراتيجية
الثانية لا تعمل إلا معها.

\end{document}
